% -*- coding: utf-8 -*-
% Aula: Caminhos minimos -- o algoritmo de Dijkstra (Unidade 3 de AED2, depois de Kruskal/Prim)
% Fontes principais: Cormen et al. (2009), cap. 24: introducao (Lema 24.1, pesos negativos,
% Initialize-Single-Source e Relax, Figura 24.3), sec. 24.3 (Dijkstra, Teorema 24.6,
% Corolario 24.7, Figura 24.6), sec. 24.5 (Lemas 24.10, 24.11, 24.14, 24.15);
% Kleinberg e Tardos (2006), sec. 4.4; Dasgupta et al. (2008), sec. 4.4-4.7;
% Levitin (2012), sec. 9.3; Feofiloff (2020).
% Exemplo condutor: o digrafo da Figura 24.6 de Cormen et al. (s,t,x,y,z; origem s).
% Dijkstra fecha s, y, z, t, x com distancias 0, 5, 7, 8, 9.
% Fechamento usa o grafo da Figura 23.1 (aula de Kruskal/Prim): arvore de caminhos minimos
% a partir de a pesa 42, a AGM pesa 37.
\documentclass[aspectratio=169,xcolor=table]{beamer}

\usetheme{UFS}

\usepackage[utf8]{inputenc}
\usepackage[T1]{fontenc}
\usepackage[brazil]{babel}
\usepackage{amsmath,amssymb,amsthm}
\usepackage{graphicx}
\usepackage{booktabs}
\usepackage{listings}
\usepackage{xcolor}
\usepackage{tikz}
\usetikzlibrary{arrows.meta, shapes, positioning, calc, fit, matrix,
  decorations.pathreplacing, decorations.pathmorphing, backgrounds, patterns}

\definecolor{codegreen}{rgb}{0,0.6,0}
\definecolor{codegray}{rgb}{0.5,0.5,0.5}
\definecolor{codepurple}{rgb}{0.58,0,0.82}
\definecolor{backcolour}{rgb}{0.95,0.95,0.92}

\lstdefinestyle{mystyle}{
    language=C,
    backgroundcolor=\color{backcolour},
    commentstyle=\color{codegreen},
    keywordstyle=\color{magenta},
    numberstyle=\tiny\color{codegray},
    stringstyle=\color{codepurple},
    breakatwhitespace=false,
    breaklines=true,
    captionpos=b,
    keepspaces=true,
    columns=fullflexible,
    numbers=none,
    showspaces=false,
    showstringspaces=false,
    showtabs=false,
    tabsize=2,
    frame=single,
    rulecolor=\color{gray},
    extendedchars=true,
    basicstyle=\ttfamily\scriptsize,
    literate={ê}{{\^e}}1 {à}{{\`a}}1 {ú}{{\'u}}1 {õ}{{\~o}}1 {ó}{{\'o}}1
             {í}{{\'i}}1 {á}{{\'a}}1 {ã}{{\~a}}1 {é}{{\'e}}1 {ç}{{\c{c}}}1
             {â}{{\^a}}1 {ô}{{\^o}}1 {ü}{{\"u}}1 {Á}{{\'A}}1 {É}{{\'E}}1
}
\lstset{style=mystyle}

% ---- estilos TikZ reutilizados ----
\tikzset{
  box/.style={rectangle, draw=blue!60, fill=blue!10, rounded corners,
              minimum width=6.6cm, minimum height=0.55cm, font=\scriptsize},
  vert/.style={circle, draw=blue!60!black, fill=blue!10, font=\tiny,
               inner sep=0pt, minimum size=4.2mm},
  vdentro/.style={vert, draw=orange!80!black, fill=orange!35},
  aresta/.style={draw=gray!55, thin},
  arvore/.style={draw=orange!85!black, line width=2.2pt},
  rejeitada/.style={draw=red!70!black, thick, dashed},
  peso/.style={font=\tiny, fill=white, inner sep=0.6pt},
  arc/.style={-{Stealth}, thick},
  arco/.style={-{Stealth[length=1.5mm]}, draw=gray!60, thin},
  arcoarv/.style={-{Stealth[length=2mm]}, draw=orange!85!black, line width=1.8pt},
  dlab/.style={font=\tiny\bfseries, text=blue!60!black, inner sep=1pt},
  r/.style={},
  bl/.style={bend left=22},
  cz/.style={out=-115, in=-80, looseness=1.25}
}

% ---- o digrafo da Figura 24.6 de Cormen et al. ----
% \digrafo{arcos destacados u/v/estilo}{vertices fechados}{d(s)}{d(t)}{d(x)}{d(y)}{d(z)}
% (usar dentro de um tikzpicture; listas vazias: passar {}; rotulo vazio: nao desenha)
\newcommand{\arcoscond}{s/t/10/r, s/y/5/r, t/x/1/r, t/y/2/bl, y/t/3/bl, y/x/9/r,
  y/z/2/r, x/z/4/bl, z/x/6/bl, z/s/7/cz}
\newcommand{\digrafo}[7]{%
  \node[vert] (s) at (0,0) {s};      \node[vert] (t) at (1.6,0.9) {t};
  \node[vert] (x) at (3.4,0.9) {x};  \node[vert] (y) at (1.6,-0.9) {y};
  \node[vert] (z) at (3.4,-0.9) {z};
  \def\lista{#2}\ifx\lista\empty\else
    \foreach \v in {#2} \node[vdentro] at (\v) {\v};\fi
  \foreach \u/\v/\w/\o in \arcoscond \draw[arco, \o] (\u) to (\v);
  \edef\lista{#1}\ifx\lista\empty\else
    \foreach \u/\v/\o in \lista \draw[arcoarv, \o] (\u) to (\v);\fi
  \foreach \u/\v/\w/\o in \arcoscond \path[\o] (\u) to node[peso] {\w} (\v);
  \def\lista{#3}\ifx\lista\empty\else \node[dlab, left=0.5mm of s] {#3};\fi
  \def\lista{#4}\ifx\lista\empty\else \node[dlab, above=0.5mm of t] {#4};\fi
  \def\lista{#5}\ifx\lista\empty\else \node[dlab, above=0.5mm of x] {#5};\fi
  \def\lista{#6}\ifx\lista\empty\else \node[dlab, below left=0mm and 0.2mm of y] {#6};\fi
  \def\lista{#7}\ifx\lista\empty\else \node[dlab, right=0.5mm of z] {#7};\fi
}
% a arvore final de caminhos minimos (pai[t]=y, pai[x]=t, pai[y]=s, pai[z]=y)
\newcommand{\arvorefinal}{s/y/r, y/t/bl, t/x/r, y/z/r}

% ---- o grafo da Figura 23.1 de Cormen et al. (aula de Kruskal e Prim) ----
\newcommand{\arestasclrs}{a/b/4, a/h/8, b/c/8, b/h/11, c/d/7, c/f/4, c/i/2,
  d/e/9, d/f/14, e/f/10, f/g/2, g/h/1, g/i/6, h/i/7}
\newcommand{\grafoclrs}[2]{%
  \coordinate (a) at (0,0);     \coordinate (b) at (1,1);
  \coordinate (c) at (2.3,1);   \coordinate (d) at (3.6,1);
  \coordinate (e) at (4.6,0);   \coordinate (f) at (3.6,-1);
  \coordinate (g) at (2.3,-1);  \coordinate (h) at (1,-1);
  \coordinate (i) at (1.75,0);
  \foreach \u/\v/\w in \arestasclrs \draw[aresta] (\u) -- (\v);
  \def\lista{#1}\ifx\lista\empty\else
    \foreach \u/\v in {#1} \draw[arvore] (\u) -- (\v);\fi
  \foreach \u/\v/\w in \arestasclrs \path (\u) -- (\v) node[peso, pos=0.5] {\w};
  \foreach \v in {a,...,i} \node[vert] at (\v) {\v};
  \def\lista{#2}\ifx\lista\empty\else
    \foreach \v in {#2} \node[vdentro] at (\v) {\v};\fi
}

% ---- ponteiro de leitura no rodape do frame ----
\newcommand{\fonte}[1]{%
  \vfill
  {\tiny\color{gray!55!black}\textbf{Para aprofundar:} #1\par}%
}

\title{Caminhos mínimos}
\subtitle{O algoritmo de Dijkstra}
\author[André Y. Kashiwabara <andreyoshiaki@dcomp.ufs.br>]{Prof.\ Dr.\ André Yoshiaki Kashiwabara}
\institute{Departamento de Computação --- Universidade Federal de Sergipe}
\date{\today}

\begin{document}

\begin{frame}[plain,noframenumbering]
  \titlepage
\end{frame}

\begin{frame}{O que você vai aprender hoje}
  \begin{center}
  \begin{tikzpicture}[node distance=2.6mm]
    \node[box] (t1) {1.\ O problema dos caminhos mínimos a partir de uma origem};
    \node[box, below=of t1] (t2) {2.\ Relaxamento: estimativas que só descem};
    \node[box, below=of t2, fill=orange!20, draw=orange!80!black] (t3)
         {3.\ Dijkstra: fechar sempre o vértice mais próximo};
    \node[box, below=of t3] (t4) {4.\ Implementação em C e custo: vetor ou heap};
    \node[box, below=of t4] (t5) {5.\ Limites: pesos negativos e o panorama};
    \foreach \a/\b in {t1/t2, t2/t3, t3/t4, t4/t5}
      \draw[-{Stealth}, thick, blue!60!black] (\a) -- (\b);
    \draw[-{Stealth}, thick, orange!80!black, dashed]
      (t3.east) to[out=0, in=0, looseness=1.4]
      node[right, xshift=1mm, font=\tiny, align=left]
      {o código de Prim\\com outra chave} (t4.east);
  \end{tikzpicture}
  \end{center}
  {\scriptsize\centering
  Na aula passada, Prim trazia para a árvore o vértice de \emph{aresta} mais barata. Hoje,
  Dijkstra traz o vértice de \emph{caminho} mais barato.\par}
  \fonte{Cormen et al.\ (2009), cap.~24~\cite{cormen2009introduction};
    Kleinberg e Tardos (2006), \S 4.4~\cite{kleinberg2006algorithm}.}
\end{frame}

% =====================================================================
\section{O problema dos caminhos mínimos}

\begin{frame}{Por que precisamos de caminhos mínimos?}
  \begin{columns}[T]
    \begin{column}{0.48\textwidth}
      \textbf{Problema:} um aplicativo de mapas precisa da rota mais rápida de Aracaju até
      cada cidade de Sergipe. As vias são de mão única ou dupla, e cada trecho tem um tempo
      de percurso.
      \medskip

      {\small A busca em largura (BFS) acha o caminho com \emph{menos trechos}, mas um
      trecho pode levar 2 minutos e outro, 40.}
    \end{column}
    \begin{column}{0.48\textwidth}
      \begin{center}
      \begin{tikzpicture}[x=0.9cm, y=0.8cm]
        \digrafo{}{}{}{}{}{}{}
      \end{tikzpicture}
      \end{center}
      {\scriptsize Exemplo condutor: origem $s$. Menos arcos: $s \to t \to x$ (custo 11).
      Menor custo: $s \to y \to t \to x$ (custo 9).\par}
    \end{column}
  \end{columns}
  \fonte{Cormen et al.\ (2009), \S 24.3, Figura 24.6~\cite{cormen2009introduction};
    Dasgupta et al.\ (2008), \S 4.4~\cite{dasgupta2008algorithms}.}
\end{frame}

\begin{frame}{Caminhos mínimos: definição}
  \begin{block}{Definição}
    Seja $G = (V, E)$ um digrafo com custos $w(u,v)$ nos arcos. O \textbf{custo} de um
    caminho $P = \langle v_0, v_1, \dots, v_k \rangle$ é
    $w(P) = \sum_{i=1}^{k} w(v_{i-1}, v_i)$. A \textbf{distância} de $s$ a $v$ é
    \[ \delta(s, v) = \min\{\, w(P) : P \text{ é caminho de } s \text{ a } v \,\}, \]
    com $\delta(s,v) = \infty$ se $v$ não é alcançável a partir de $s$.
  \end{block}
  \begin{exampleblock}{Intuição}
    \small Problema de \textbf{origem única}: uma origem $s$ fixa, a distância a \emph{todos}
    os vértices de uma vez. Achar só $s \to v$ não é mais fácil no pior caso.
  \end{exampleblock}
  {\scriptsize Notação: $n = |V|$, $m = |E|$. Grafo não dirigido: cada aresta vira dois arcos.
  Hoje, todos os custos são $\ge 0$.\par}
  \fonte{Cormen et al.\ (2009), introdução do cap.~24~\cite{cormen2009introduction};
    Kleinberg e Tardos (2006), \S 4.4~\cite{kleinberg2006algorithm}.}
\end{frame}

\begin{frame}{Exemplo: as distâncias no digrafo condutor}
  \begin{columns}[c]
    \begin{column}{0.52\textwidth}
      \begin{center}
      \begin{tikzpicture}[x=1.15cm, y=1.05cm]
        \digrafo{\arvorefinal}{}{0}{8}{9}{5}{7}
      \end{tikzpicture}
      \end{center}
    \end{column}
    \begin{column}{0.44\textwidth}
      \centering\scriptsize
      \begin{tabular}{clc}
        \toprule
        $v$ & caminho mínimo & $\delta(s,v)$ \\
        \midrule
        $s$ & $s$ & 0 \\
        $y$ & $s \to y$ & 5 \\
        $z$ & $s \to y \to z$ & 7 \\
        $t$ & $s \to y \to t$ & 8 \\
        $x$ & $s \to y \to t \to x$ & 9 \\
        \bottomrule
      \end{tabular}
      \medskip

      \raggedright
      Em azul, $\delta(s,v)$. Os arcos em laranja formam uma \textbf{árvore de caminhos
      mínimos}: o caminho de $s$ a cada $v$ dentro dela é mínimo. Note que o arco
      $s \to t$, de custo 10, não serve para nada.
    \end{column}
  \end{columns}
  \fonte{Cormen et al.\ (2009), \S 24.3, Figura 24.6(f)~\cite{cormen2009introduction}.}
\end{frame}

\begin{frame}{A subestrutura ótima dos caminhos mínimos}
  \begin{columns}[c]
    \begin{column}{0.55\textwidth}
      \begin{block}{Lema}
        \small Todo trecho de um caminho mínimo é um caminho mínimo.
      \end{block}
      \scriptsize
      \textbf{Prova (recorte e cola).} Seja $P$ mínimo de $s$ a $v$ e $P_{ab}$ o trecho de $a$
      a $b$. Se houvesse $Q$ de $a$ a $b$ com $w(Q) < w(P_{ab})$, trocar $P_{ab}$ por $Q$ em
      $P$ daria um caminho de $s$ a $v$ mais barato que $P$: absurdo. \hfill$\square$
      \medskip

      \textbf{Consequência.} Basta guardar, para cada $v$, o \textbf{predecessor}
      $\mathit{pai}[v]$ no caminho mínimo: os $n - 1$ arcos $(\mathit{pai}[v], v)$ formam a
      árvore da página anterior.
    \end{column}
    \begin{column}{0.41\textwidth}
      \begin{tikzpicture}[x=0.75cm, y=0.6cm, p/.style={circle, fill=blue!60!black,
          inner sep=0pt, minimum size=1.6mm}]
        \node[p, label={[font=\tiny]below:$s$}] (s) at (0,0) {};
        \node[p, label={[font=\tiny]below:$a$}] (a) at (1.6,0) {};
        \node[p, label={[font=\tiny]below:$b$}] (b) at (3.4,0) {};
        \node[p, label={[font=\tiny]below:$v$}] (v) at (5,0) {};
        \draw[thick, blue!60!black, decorate, decoration={snake, amplitude=0.6mm,
          segment length=3mm}] (s) -- (a);
        \draw[very thick, red!70!black, decorate, decoration={snake, amplitude=0.6mm,
          segment length=3mm}] (a) -- node[below=1mm, font=\tiny] {$P_{ab}$} (b);
        \draw[thick, blue!60!black, decorate, decoration={snake, amplitude=0.6mm,
          segment length=3mm}] (b) -- (v);
        \draw[very thick, orange!85!black, dashed] (a) to[bend left=55]
          node[above, font=\tiny] {$Q$ mais barato?} (b);
      \end{tikzpicture}
      \smallskip

      \tiny É a \textbf{subestrutura ótima} das aulas de programação dinâmica; Floyd já
      dependia dela.
    \end{column}
  \end{columns}
  \fonte{Cormen et al.\ (2009), cap.~24, Lema 24.1~\cite{cormen2009introduction};
    Kleinberg e Tardos (2006), \S 4.4~\cite{kleinberg2006algorithm}.}
\end{frame}

\begin{frame}{O que já sabemos fazer}
  \begin{center}
  \scriptsize
  \begin{tabular}{llll}
    \toprule
    algoritmo & resolve & restrição & tempo \\
    \midrule
    força bruta & enumera todos os caminhos & --- & exponencial \\
    BFS (Unidade 1) & origem única & todo custo $= 1$ & $O(n + m)$ \\
    Floyd (Unidade 2) & \emph{todos os pares} & sem ciclo negativo & $\Theta(n^3)$ \\
    \textbf{Dijkstra (hoje)} & origem única & custos $\ge 0$ & $O(m \log n)$ \\
    \bottomrule
  \end{tabular}
  \end{center}
  \medskip
  \begin{columns}[T]
    \begin{column}{0.48\textwidth}
      \scriptsize
      Floyd responde à pergunta de origem única, mas paga $\Theta(n^3)$ calculando $n^2$
      distâncias quando queremos só $n$.
    \end{column}
    \begin{column}{0.48\textwidth}
      \scriptsize
      Num mapa de ruas com $n = 10^6$ cruzamentos, $n^3 = 10^{18}$ operações; o grafo é
      esparso ($m \approx 3n$), e $m \log n \approx 6 \times 10^7$.
    \end{column}
  \end{columns}
  \fonte{Cormen et al.\ (2009), cap.~24 e \S 25.2~\cite{cormen2009introduction};
    Dasgupta et al.\ (2008), \S 4.4~\cite{dasgupta2008algorithms}.}
\end{frame}

\begin{frame}{Recapitulando: o problema}
  \begin{block}{O que aprendemos}
    \begin{itemize}
      \item $\delta(s,v)$ é o menor custo de um caminho de $s$ a $v$ ($\infty$ se não há).
      \item Origem única: queremos $\delta(s,v)$ para todo $v$ de uma vez.
      \item Trechos de caminhos mínimos são mínimos: basta guardar $\mathit{pai}[v]$.
      \item Os predecessores formam uma árvore de caminhos mínimos.
      \item BFS só vale para custos unitários; Floyd é caro demais para uma origem.
    \end{itemize}
  \end{block}
  \fonte{Cormen et al.\ (2009), cap.~24~\cite{cormen2009introduction}.}
\end{frame}

% =====================================================================
\section{Relaxamento}

\begin{frame}{Por que trabalhar com estimativas?}
  \begin{columns}[T]
    \begin{column}{0.48\textwidth}
      \textbf{Problema:} não conhecemos $\delta(s,v)$ de antemão. Precisamos de uma operação
      que melhore aos poucos o que sabemos, sem nunca errar para baixo.
      \medskip

      {\small Analogia: o preço de uma passagem que você vai pesquisando. Cada rota nova
      encontrada só pode \emph{baixar} o melhor preço conhecido.}
    \end{column}
    \begin{column}{0.48\textwidth}
      \textbf{Solução:} para cada $v$, uma \textbf{estimativa} $\mathit{dist}[v]$ que é o custo
      de \emph{algum} caminho de $s$ a $v$ já encontrado:
      \begin{center}
      \fbox{\small $\mathit{dist}[v] \ge \delta(s,v)$ sempre}
      \end{center}
      Começa com $\mathit{dist}[s] = 0$ e $\infty$ no resto, e só desce.
      \medskip

      {\small Dijkstra, Bellman-Ford e o algoritmo para DAGs usam a mesma operação; o que muda
      é a \emph{ordem} em que a aplicam.}
    \end{column}
  \end{columns}
  \fonte{Cormen et al.\ (2009), cap.~24, ``Relaxation''~\cite{cormen2009introduction}.}
\end{frame}

\begin{frame}{Relaxar um arco: definição}
  \begin{columns}[c]
    \begin{column}{0.5\textwidth}
      \begin{block}{\textsc{Relaxa}$(u, v)$}
        \small
        \textbf{se} $\mathit{dist}[u] + w(u,v) < \mathit{dist}[v]$:\\
        \quad $\mathit{dist}[v] \gets \mathit{dist}[u] + w(u,v)$\\
        \quad $\mathit{pai}[v] \gets u$
      \end{block}
      \begin{exampleblock}{Intuição}
        \scriptsize ``Chegar a $v$ passando por $u$ é melhor do que o melhor caminho que eu
        conhecia?'' Se sim, adote o novo caminho: o de $u$ seguido do arco $u \to v$.
      \end{exampleblock}
    \end{column}
    \begin{column}{0.46\textwidth}
      \begin{center}
      \begin{tikzpicture}[x=1cm, y=0.9cm]
        \foreach \k/\yy/\du/\dv/\dn/\cor in {1/0/5/9/7/orange!85!black, 2/-1.6/5/6/6/gray} {
          \node[vert] (u\k) at (0,\yy) {u};
          \node[vert] (v\k) at (1.5,\yy) {v};
          \draw[arco, draw=gray!70] (u\k) -- node[peso, above] {2} (v\k);
          \node[dlab, below=0.5mm of u\k] {\du};
          \node[dlab, below=0.5mm of v\k] {\dv};
          \draw[-{Stealth}, thick, gray] (2.0,\yy) -- (2.7,\yy);
          \node[vert] (U\k) at (3.2,\yy) {u};
          \node[vert, draw=\cor] (V\k) at (4.7,\yy) {v};
          \draw[arco, draw=\cor, thick] (U\k) -- node[peso, above] {2} (V\k);
          \node[dlab, below=0.5mm of U\k] {\du};
          \node[dlab, below=0.5mm of V\k] {\dn};
        }
      \end{tikzpicture}
      \end{center}
      \scriptsize
      Em cima, $5 + 2 < 9$: $\mathit{dist}[v]$ cai para 7 e $\mathit{pai}[v] = u$. Embaixo,
      $5 + 2 \ge 6$: nada muda.
    \end{column}
  \end{columns}
  \fonte{Cormen et al.\ (2009), cap.~24, \textsc{Relax} e Figura 24.3~\cite{cormen2009introduction}.}
\end{frame}

\begin{frame}{Propriedades do relaxamento}
  \begin{columns}[T]
    \begin{column}{0.5\textwidth}
      \scriptsize
      \textbf{Desigualdade triangular.} Para todo arco $(u,v)$:
      $\delta(s,v) \le \delta(s,u) + w(u,v)$.
      \medskip

      \textbf{Limite superior.} Em qualquer sequência de relaxamentos,
      $\mathit{dist}[v] \ge \delta(s,v)$, e quando $\mathit{dist}[v]$ atinge $\delta(s,v)$ não
      muda mais.
      \medskip

      \textbf{Convergência.} Se $s \leadsto u \to v$ é um caminho mínimo e
      $\mathit{dist}[u] = \delta(s,u)$ no momento em que $(u,v)$ é relaxado, então depois
      disso $\mathit{dist}[v] = \delta(s,v)$.
    \end{column}
    \begin{column}{0.46\textwidth}
      \scriptsize
      \textbf{Relaxamento de caminho.} Se os arcos de um caminho mínimo
      $\langle s, v_1, \dots, v_k \rangle$ forem relaxados \emph{nessa ordem} (com quaisquer
      outros relaxamentos no meio), no fim $\mathit{dist}[v_k] = \delta(s, v_k)$.
      \medskip

      \begin{block}{A pergunta da aula}
        \scriptsize Em que ordem relaxar os arcos para que cada arco seja relaxado
        \textbf{uma vez só}?
      \end{block}
      \medskip

      Para custos $\ge 0$, a resposta é de Dijkstra.
    \end{column}
  \end{columns}
  \fonte{Cormen et al.\ (2009), \S 24.5, Lemas 24.10, 24.11, 24.14 e
    24.15~\cite{cormen2009introduction}.}
\end{frame}

\begin{frame}{Recapitulando: relaxamento}
  \begin{block}{O que aprendemos}
    \begin{itemize}
      \item $\mathit{dist}[v]$ é o custo de um caminho conhecido: nunca menor que $\delta(s,v)$.
      \item Relaxar $(u,v)$ testa se passar por $u$ melhora o caminho até $v$.
      \item Relaxar os arcos de um caminho mínimo em ordem acerta a distância do fim.
      \item Os algoritmos de caminho mínimo diferem na ordem dos relaxamentos.
    \end{itemize}
  \end{block}
  \fonte{Cormen et al.\ (2009), cap.~24 e \S 24.5~\cite{cormen2009introduction}.}
\end{frame}

% =====================================================================
\section{O algoritmo de Dijkstra}

\begin{frame}{Por que fechar o vértice mais próximo?}
  \begin{columns}[T]
    \begin{column}{0.48\textwidth}
      \textbf{Problema:} relaxar um arco $(u,v)$ só é definitivo se $\mathit{dist}[u]$ já for
      exata. Qual vértice tem certamente a estimativa exata?
      \medskip

      {\small Imagine a origem soltando água que escorre pelos arcos com velocidade
      constante: os vértices molham em ordem crescente de distância. O mais próximo dos ainda
      secos é o próximo a molhar.}
    \end{column}
    \begin{column}{0.48\textwidth}
      \textbf{Solução (Dijkstra, 1959):} manter o conjunto $S$ de vértices
      \textbf{fechados}, cuja distância já é exata.
      \begin{enumerate}
        \small
        \item $S \gets \emptyset$; $\mathit{dist}[s] = 0$;
        \item escolha $u \notin S$ com \textbf{menor} $\mathit{dist}[u]$ e feche $u$;
        \item relaxe os arcos que saem de $u$;
        \item repita até fechar todos os alcançáveis.
      \end{enumerate}
      {\small Cada arco é relaxado uma vez só: quando sua ponta inicial fecha.}
    \end{column}
  \end{columns}
  \fonte{Dijkstra (1959)~\cite{dijkstra1959note}; Dasgupta et al.\ (2008), \S 4.4
    (adaptação da BFS)~\cite{dasgupta2008algorithms}; Levitin (2012),
    \S 9.3~\cite{levitin2012introduction}.}
\end{frame}

\begin{frame}{O algoritmo de Dijkstra}
  \begin{columns}[c]
    \begin{column}{0.48\textwidth}
      \begin{block}{\textsc{Dijkstra}$(G, w, s)$}
        \small
        \textbf{para cada} $v$: $\mathit{dist}[v] \gets \infty$; $\mathit{pai}[v] \gets -1$\\
        $\mathit{dist}[s] \gets 0$;\ $S \gets \emptyset$\\
        $Q \gets$ fila de prioridade com $V$, chave $\mathit{dist}$\\
        \textbf{enquanto} $Q \neq \emptyset$:\\
        \quad $u \gets$ \textsc{Extrai-Min}$(Q)$;\ $S \gets S \cup \{u\}$\\
        \quad \textbf{para cada} arco $(u,v)$:\\
        \quad\quad \textsc{Relaxa}$(u,v)$
      \end{block}
    \end{column}
    \begin{column}{0.48\textwidth}
      \scriptsize
      Compare com Prim, da aula passada: a estrutura é a mesma, só a chave muda.
      \medskip

      \begin{tabular}{lp{3.6cm}}
        \toprule
        & a chave de $v$ baixa para \\
        \midrule
        \textbf{Prim} & $w(u,v)$ \\
        \addlinespace
        \textbf{Dijkstra} & $\mathit{dist}[u] + w(u,v)$ \\
        \bottomrule
      \end{tabular}
      \medskip

      Prim pergunta ``quanto custa \emph{pendurar} $v$ na árvore?''; Dijkstra pergunta
      ``quanto custa \emph{chegar} a $v$ desde $s$?''. Com custos unitários e uma fila comum
      no lugar da de prioridade, Dijkstra vira a BFS.
    \end{column}
  \end{columns}
  \fonte{Cormen et al.\ (2009), \S 24.3, \textsc{Dijkstra}~\cite{cormen2009introduction};
    Kleinberg e Tardos (2006), \S 4.4~\cite{kleinberg2006algorithm}.}
\end{frame}

\begin{frame}{Exemplo: Dijkstra a partir de $s$}
  \begin{columns}[c]
    \begin{column}{0.5\textwidth}
      \centering\scriptsize
      \setlength{\tabcolsep}{3.5pt}
      \begin{tabular}{c|ccccc}
        \toprule
        fecha & $s$ & $t$ & $x$ & $y$ & $z$ \\
        \midrule
        --- & 0 & $\infty$ & $\infty$ & $\infty$ & $\infty$ \\
        $s$ (0) & -- & \textbf{10} & $\infty$ & \textbf{5} & $\infty$ \\
        $y$ (5) & -- & \textbf{8} & \textbf{14} & -- & \textbf{7} \\
        $z$ (7) & -- & 8 & \textbf{13} & -- & -- \\
        $t$ (8) & -- & -- & \textbf{9} & -- & -- \\
        $x$ (9) & -- & -- & -- & -- & -- \\
        \bottomrule
      \end{tabular}
      \smallskip

      \raggedright
      $\mathit{dist}$ depois de cada fechamento; em negrito, as que baixaram.
    \end{column}
    \begin{column}{0.46\textwidth}
      \scriptsize
      \textbf{Fecha $y$:} $t$ cai de 10 para $5 + 3 = 8$, pois chegar a $t$ por $y$ é mais
      barato que pelo arco direto.
      \medskip

      \textbf{Fecha $z$:} $x$ cai de 14 para $7 + 6 = 13$. O arco $z \to s$ não muda nada:
      $s$ já está fechado.
      \medskip

      \textbf{Fecha $t$:} $x$ cai de novo, para $8 + 1 = 9$. Ao fechar $x$, todos os
      relaxamentos falham.
      \medskip

      Ordem de fechamento: $s, y, z, t, x$, com distâncias \textbf{crescentes}
      $0, 5, 7, 8, 9$.
    \end{column}
  \end{columns}
  \fonte{Cormen et al.\ (2009), \S 24.3, Figura 24.6~\cite{cormen2009introduction}.}
\end{frame}

\begin{frame}{As estimativas ao longo do tempo}
  \begin{center}
  \begin{tikzpicture}[x=0.82cm, y=0.7cm]
    \begin{scope}
      \digrafo{s/y/r, s/t/r}{s}{0}{10}{$\infty$}{5}{$\infty$}
      \node[font=\tiny] at (1.7,-2.55) {depois de fechar $s$};
    \end{scope}
    \begin{scope}[xshift=5.2cm]
      \digrafo{s/y/r, y/t/bl, y/x/r, y/z/r}{s,y}{0}{8}{14}{5}{7}
      \node[font=\tiny] at (1.7,-2.55) {depois de fechar $y$};
    \end{scope}
    \begin{scope}[xshift=10.4cm]
      \digrafo{\arvorefinal}{s,y,z,t}{0}{8}{9}{5}{7}
      \node[font=\tiny] at (1.7,-2.55) {depois de fechar $t$};
    \end{scope}
  \end{tikzpicture}
  \end{center}
  {\small Em laranja, os vértices fechados e os arcos $(\mathit{pai}[v], v)$ do momento. O
  arco $s \to t$ entra na árvore e depois sai, quando $y$ oferece um caminho melhor até $t$; o
  mesmo acontece com $y \to x$.\par}
  \fonte{Cormen et al.\ (2009), \S 24.3, Figura 24.6(b)--(e)~\cite{cormen2009introduction}.}
\end{frame}

\begin{frame}{Dijkstra: por que está correto}
  \begin{block}{Teorema}
    \small Com custos $w \ge 0$, ao ser fechado, cada vértice $u$ tem
    $\mathit{dist}[u] = \delta(s,u)$.
  \end{block}
  \begin{columns}[c]
    \begin{column}{0.56\textwidth}
      \scriptsize
      \textbf{Prova.} Suponha que $u$ seja o primeiro a fechar com
      $\mathit{dist}[u] > \delta(s,u)$. Tome um caminho mínimo $P$ de $s$ a $u$ e seja $y$ o
      primeiro vértice de $P$ fora de $S$, com predecessor $x \in S$.
      \smallskip

      $x$ fechou com $\mathit{dist}[x] = \delta(s,x)$ e relaxou $(x,y)$; por convergência,
      $\mathit{dist}[y] = \delta(s,y)$. Como $y$ vem antes de $u$ em $P$ e os custos são
      $\ge 0$:
      \[ \mathit{dist}[y] = \delta(s,y) \le \delta(s,u) < \mathit{dist}[u] . \]
      Mas o algoritmo escolheu $u$, e não $y$: $\mathit{dist}[u] \le \mathit{dist}[y]$.
      Absurdo. \hfill$\square$
    \end{column}
    \begin{column}{0.4\textwidth}
      \begin{tikzpicture}[x=0.72cm, y=0.6cm, p/.style={circle, fill=blue!60!black,
          inner sep=0pt, minimum size=1.6mm}]
        \fill[orange!12] (-0.4,-1.4) rectangle (2.5,1.4);
        \node[font=\tiny, anchor=north west] at (-0.4,1.4) {$S$ (fechados)};
        \node[p, label={[font=\tiny]left:$s$}] (s) at (0.2,-0.6) {};
        \node[p, label={[font=\tiny]above:$x$}] (x) at (1.9,0.5) {};
        \node[p, label={[font=\tiny]above:$y$}] (y) at (3.4,0.5) {};
        \node[p, label={[font=\tiny]right:$u$}] (u) at (5.0,-0.6) {};
        \draw[thick, blue!60!black, decorate, decoration={snake, amplitude=0.5mm,
          segment length=2.5mm}] (s) -- (x);
        \draw[-{Stealth}, very thick, orange!85!black] (x) -- (y);
        \draw[thick, blue!60!black, decorate, decoration={snake, amplitude=0.5mm,
          segment length=2.5mm}] (y) -- (u);
        \draw[-{Stealth}, gray, dashed] (s) to[bend right=25]
          node[below, font=\tiny] {escolhido} (u);
      \end{tikzpicture}
      \smallskip

      \tiny Em azul e laranja, o caminho mínimo $P$. É aqui que $w \ge 0$ é indispensável:
      o trecho de $y$ a $u$ não pode diminuir o custo.
    \end{column}
  \end{columns}
  \fonte{Cormen et al.\ (2009), \S 24.3, Teorema 24.6 e Figura 24.7~\cite{cormen2009introduction};
    Kleinberg e Tardos (2006), \S 4.4, (4.14)~\cite{kleinberg2006algorithm}.}
\end{frame}

\begin{frame}{Recapitulando: Dijkstra}
  \begin{block}{O que aprendemos}
    \begin{itemize}
      \item Dijkstra fecha sempre o vértice aberto de menor $\mathit{dist}$.
      \item Ao fechar $u$, relaxa os arcos que saem de $u$; cada arco, uma vez.
      \item É o código de Prim com chave $\mathit{dist}[u] + w(u,v)$ no lugar de $w(u,v)$.
      \item Os vértices fecham em ordem crescente de distância.
      \item A correção depende de $w \ge 0$: nenhum caminho fica mais barato ao crescer.
    \end{itemize}
  \end{block}
  \fonte{Cormen et al.\ (2009), \S 24.3~\cite{cormen2009introduction}.}
\end{frame}

% =====================================================================
\section{Implementação e custo}

\begin{frame}[fragile]{Dijkstra em C com matriz de adjacência}
  \begin{columns}[T]
    \begin{column}{0.58\textwidth}
\begin{lstlisting}[basicstyle=\ttfamily\tiny]
/* W[u][v] = custo do arco u->v, ou INF se nao ha arco */
void dijkstra( int n, int W[][MAXN], int s, int dist[], int pai[]) {
   bool fechado[MAXN];
   for (int v = 0; v < n; ++v) {
      dist[v] = INF;  pai[v] = -1;  fechado[v] = false;
   }
   dist[s] = 0;
   for (int it = 0; it < n; ++it) {
      int u = -1;                     /* o aberto mais proximo */
      for (int v = 0; v < n; ++v)
         if (!fechado[v] && (u < 0 || dist[v] < dist[u]))
            u = v;
      if (dist[u] == INF) break;      /* o resto e inalcancavel */
      fechado[u] = true;
      for (int v = 0; v < n; ++v)     /* relaxa os arcos de u */
         if (W[u][v] != INF && dist[u] + W[u][v] < dist[v]) {
            dist[v] = dist[u] + W[u][v];  pai[v] = u;
         }
   }
}
\end{lstlisting}
    \end{column}
    \begin{column}{0.38\textwidth}
      \scriptsize
      \textbf{Diferenças para o Prim da aula passada:}
      \begin{itemize}
        \scriptsize
        \item a chave é $\mathit{dist}[u] + W[u][v]$;
        \item o teste \texttt{W[u][v] != INF} evita somar com \texttt{INF} e transbordar o
              \texttt{int};
        \item se o mais próximo tem $\mathit{dist} = \infty$, os demais são inalcançáveis.
      \end{itemize}
      \medskip

      \textbf{Custo.} $n$ iterações, cada uma com duas varreduras de $n$ posições:
      $\Theta(n^2)$, ótimo para grafos densos.
    \end{column}
  \end{columns}
  \fonte{Feofiloff (2020), algoritmo de Dijkstra~\cite{feofiloff2020grafos};
    Levitin (2012), \S 9.3~\cite{levitin2012introduction}.}
\end{frame}

\begin{frame}[fragile]{Reconstruindo o caminho com \texttt{pai[]}}
  \begin{columns}[T]
    \begin{column}{0.5\textwidth}
\begin{lstlisting}[basicstyle=\ttfamily\tiny]
/* imprime o caminho minimo de s a t, de s para t */
void imprime_caminho( int pai[], int s, int t) {
   if (t == s) { printf( "%d", s);  return; }
   if (pai[t] == -1) {
      printf( "%d inalcancavel", t);  return;
   }
   imprime_caminho( pai, s, pai[t]);
   printf( " -> %d", t);
}
\end{lstlisting}
      \scriptsize
      A recursão sobe de $t$ até $s$ e imprime na volta, na ordem certa. Custo: o número de
      arcos do caminho.
    \end{column}
    \begin{column}{0.46\textwidth}
      \begin{center}
      \scriptsize
      \begin{tabular}{cccccc}
        \toprule
        $v$ & $s$ & $t$ & $x$ & $y$ & $z$ \\
        \midrule
        $\mathit{pai}[v]$ & $-1$ & $y$ & $t$ & $s$ & $y$ \\
        \bottomrule
      \end{tabular}
      \end{center}
      \scriptsize
      Caminho até $x$: $\mathit{pai}[x] = t$, $\mathit{pai}[t] = y$, $\mathit{pai}[y] = s$.
      Lido de trás para frente: $s \to y \to t \to x$, custo $5 + 3 + 1 = 9$.
      \medskip

      O vetor $\mathit{pai}$ guarda a árvore de caminhos mínimos inteira em $n$ posições:
      os $n$ caminhos juntos podem ter $\Theta(n^2)$ arcos.
    \end{column}
  \end{columns}
  \fonte{Cormen et al.\ (2009), \S 22.2, \textsc{Print-Path}, e cap.~24, ``Representing
    shortest paths''~\cite{cormen2009introduction}.}
\end{frame}

\begin{frame}{Quanto custa Dijkstra?}
  \begin{center}
  \scriptsize
  \begin{tabular}{lll}
    \toprule
    fila de prioridade & operações dominantes & total \\
    \midrule
    vetor (varredura) & $n$ extrações de $\Theta(n)$ & $\Theta(n^2)$ \\
    heap binário & $n$ extrações e até $m$ reduções de chave, $O(\log n)$ cada & $O((n + m) \log n)$ \\
    heap de Fibonacci & redução de chave em $O(1)$ amortizado & $O(m + n \log n)$ \\
    \bottomrule
  \end{tabular}
  \end{center}
  \medskip
  \begin{columns}[T]
    \begin{column}{0.48\textwidth}
      \scriptsize
      \textbf{Grafo esparso} ($m = O(n)$): heap binário dá $O(n \log n)$; o vetor perde com
      $\Theta(n^2)$.
      \medskip

      \textbf{Grafo denso} ($m = \Theta(n^2)$): o vetor dá $\Theta(n^2)$; o heap binário
      perde com $O(n^2 \log n)$.
    \end{column}
    \begin{column}{0.48\textwidth}
      \scriptsize
      \textbf{Heap preguiçoso.} Sem redução de chave: a cada relaxamento bem-sucedido, insira
      o par $(\mathit{dist}[v], v)$ de novo; ao extrair um par com custo maior que
      $\mathit{dist}[v]$, descarte-o. O heap chega a $m$ pares, e o custo fica
      $O(m \log n)$. É o Exercício 2 do tutorial.
    \end{column}
  \end{columns}
  \fonte{Cormen et al.\ (2009), \S 24.3 (análise)~\cite{cormen2009introduction};
    Fredman e Tarjan (1987)~\cite{fredman1987fibonacci}; Dasgupta et al.\ (2008),
    \S 4.5~\cite{dasgupta2008algorithms}.}
\end{frame}

\begin{frame}{Recapitulando: implementação}
  \begin{block}{O que aprendemos}
    \begin{itemize}
      \item Com matriz e varredura, Dijkstra custa $\Theta(n^2)$: bom para grafos densos.
      \item Com listas de adjacência e heap binário, $O((n + m) \log n)$: bom para esparsos.
      \item O heap preguiçoso dispensa a redução de chave e descarta pares velhos.
      \item O vetor $\mathit{pai}$ guarda todos os caminhos em $n$ posições.
    \end{itemize}
  \end{block}
  \fonte{Cormen et al.\ (2009), \S 24.3~\cite{cormen2009introduction}.}
\end{frame}

% =====================================================================
\section{Limites e panorama}

\begin{frame}{Por que custos negativos quebram Dijkstra?}
  \begin{columns}[c]
    \begin{column}{0.4\textwidth}
      \begin{center}
      \begin{tikzpicture}[x=1.1cm, y=1.0cm]
        \node[vert] (s) at (0,0) {s};
        \node[vdentro] (a) at (1.5,0.8) {a};
        \node[vdentro] (c) at (3,0.8) {c};
        \node[vert] (b) at (1.5,-0.8) {b};
        \draw[arco, draw=gray!80] (s) -- node[peso] {2} (a);
        \draw[arco, draw=gray!80] (s) -- node[peso] {4} (b);
        \draw[arco, draw=gray!80] (a) -- node[peso] {1} (c);
        \draw[arco, draw=red!70!black, thick] (b) -- node[peso] {$-3$} (a);
        \node[dlab, above=0.5mm of a] {2};
        \node[dlab, above=0.5mm of c] {3};
        \node[dlab, below=0.5mm of b] {4};
      \end{tikzpicture}
      \end{center}
    \end{column}
    \begin{column}{0.56\textwidth}
      \scriptsize
      Dijkstra fecha $s$ (0), $a$ (2), $c$ (3) e por último $b$ (4). Ao fechar $b$, o arco
      $b \to a$ ofereceria $4 - 3 = 1$ para $a$, mas $a$ e $c$ já estão fechados.
      \medskip

      \textbf{Resposta errada:} $\mathit{dist}[c] = 3$, mas
      $s \to b \to a \to c$ custa $4 - 3 + 1 = 2$.
      \medskip

      A prova falha exatamente no passo $\delta(s,y) \le \delta(s,u)$: um trecho com custo
      negativo pode deixar o caminho mais barato ao crescer.
      \medskip

      \textbf{E somar uma constante a todos os custos?} Não funciona: penaliza mais os
      caminhos com mais arcos. Somando 3, $s \to a \to c$ custa 9 e $s \to b \to a \to c$
      custa 11, e o mínimo muda.
    \end{column}
  \end{columns}
  \fonte{Cormen et al.\ (2009), cap.~24, ``Negative-weight edges'', e Exercício
    24.3-2~\cite{cormen2009introduction}; Dasgupta et al.\ (2008),
    \S 4.6~\cite{dasgupta2008algorithms}.}
\end{frame}

\begin{frame}{Árvore de caminhos mínimos não é AGM}
  \begin{columns}[c]
    \begin{column}{0.48\textwidth}
      \begin{center}
      \begin{tikzpicture}[x=0.95cm, y=0.95cm]
        \grafoclrs{g/h, c/i, f/g, a/b, c/f, c/d, a/h, d/e}{a}
      \end{tikzpicture}
      \end{center}
      \scriptsize\centering AGM de Kruskal: peso \textbf{37}.
      De $a$ a $e$: $a\,h\,g\,f\,c\,d\,e$, custo 31.
    \end{column}
    \begin{column}{0.48\textwidth}
      \begin{center}
      \begin{tikzpicture}[x=0.95cm, y=0.95cm]
        \grafoclrs{a/b, b/c, c/d, f/e, g/f, h/g, a/h, c/i}{a}
      \end{tikzpicture}
      \end{center}
      \scriptsize\centering Dijkstra a partir de $a$: peso \textbf{42}.
      De $a$ a $e$: $a\,h\,g\,f\,e$, custo 21.
    \end{column}
  \end{columns}
  \medskip
  {\scriptsize O grafo da aula passada. A AGM minimiza a \textbf{soma de todas as arestas}; a
  árvore de Dijkstra minimiza \textbf{cada caminho a partir de $a$}. Os objetivos são
  diferentes e as árvores também: a de Dijkstra troca $(c,f)$ e $(d,e)$ por $(b,c)$ e
  $(e,f)$.\par}
  \fonte{Cormen et al.\ (2009), \S 23.1, Figura 23.1, e \S 24.3~\cite{cormen2009introduction}.}
\end{frame}

\begin{frame}{O panorama dos caminhos mínimos}
  \begin{center}
  \scriptsize
  \begin{tabular}{llll}
    \toprule
    algoritmo & quando usar & ordem dos relaxamentos & tempo \\
    \midrule
    BFS & custos unitários & fila comum & $O(n + m)$ \\
    DAG & digrafo acíclico, qualquer custo & ordem topológica & $O(n + m)$ \\
    \textbf{Dijkstra} & custos $\ge 0$ & menor $\mathit{dist}$ primeiro & $O(m \log n)$ \\
    Bellman-Ford & custos negativos; detecta ciclo negativo & todos os arcos, $n - 1$ vezes & $O(nm)$ \\
    Floyd & todos os pares & intermediários $0..k-1$ & $\Theta(n^3)$ \\
    Johnson & todos os pares, esparso & Bellman-Ford + $n \times$ Dijkstra & $O(nm \log n)$ \\
    \bottomrule
  \end{tabular}
  \end{center}
  \medskip
  {\scriptsize Todos, exceto Floyd, são o mesmo relaxamento em ordens diferentes. A ordem
  topológica da Unidade 1 também resolve caminhos mínimos, e até \emph{máximos}, em DAGs.
  Na prática, GPS e roteadores usam Dijkstra e variantes com heurística, como o A*
  (desafio do tutorial).\par}
  \fonte{Cormen et al.\ (2009), \S 24.1, \S 24.2, \S 24.3 e \S 25.3~\cite{cormen2009introduction};
    Hart, Nilsson e Raphael (1968), A*~\cite{hart1968formal}.}
\end{frame}

\begin{frame}{Fechando}
  \begin{columns}[T]
    \begin{column}{0.48\textwidth}
      \begin{block}{O que levar desta aula}
        \scriptsize
        \begin{itemize}
          \item Relaxar arcos só baixa estimativas, que nunca ficam abaixo de $\delta$.
          \item Dijkstra relaxa a partir do vértice aberto mais próximo, uma vez por arco.
          \item É guloso: fecha cada vértice sem voltar atrás, como Prim.
          \item Exige custos $\ge 0$; para negativos, Bellman-Ford.
          \item $\Theta(n^2)$ com vetor, $O(m \log n)$ com heap.
        \end{itemize}
      \end{block}
    \end{column}
    \begin{column}{0.48\textwidth}
      \scriptsize
      \textbf{História.} Dijkstra conta que concebeu o algoritmo em 1956, em uns vinte
      minutos, sentado num café em Amsterdã, sem papel e lápis, para demonstrar um computador
      novo calculando rotas entre 64 cidades holandesas. Publicou três anos depois, numa nota
      de três páginas que também resolvia a AGM.
      \medskip

      \textbf{Próxima aula:} fluxo máximo e emparelhamento, em que caminhos num grafo
      \emph{residual} são usados para aumentar o fluxo.
    \end{column}
  \end{columns}
  \fonte{Misa e Frana (2010)~\cite{misa2010interview}; Dijkstra (1959)~\cite{dijkstra1959note}.}
\end{frame}

% =====================================================================
\section*{Referências}
\begin{frame}[allowframebreaks]{Referências}
  \tiny
  \nocite{cormen2009introduction,cormen2012algoritmos,kleinberg2006algorithm,dasgupta2008algorithms,levitin2012introduction,feofiloff2020grafos,sedgewick2011algorithms,dijkstra1959note,fredman1987fibonacci,hart1968formal,misa2010interview}
  \bibliographystyle{plain}
  \bibliography{../referencias}
\end{frame}

\end{document}
